\section{Práctica y retos del Web Agent}
\subsection{Planear un viaje a la GTC con un Web Agent}
En la demostración de 00:15:41 a 00:16:39, el Agent primero le pregunta al usuario «quiero ir a GTC 2024», después visita sucesivamente Google y Google Maps, y finalmente encuentra San Jose Convention Center; luego deja que el navegador ejecute el formulario y la navegación de la ruta. Este fragmento muestra que el Web Agent puede completar la tarea apoyándose en la interacción con la página incluso sin disponer de una API, y también muestra que necesita preguntarle constantemente al ser humano y hacer replanning dinámico.
\footnote{El intervalo de video correspondiente a este fragmento es 00:15:41--00:16:39.}

\begin{importantbox}{Los límites de la colaboración entre el humano y el Web Agent}
En los escenarios donde no se dispone de una API, el Web Agent depende sobre todo de la percepción visual y de un ciclo de instrucciones humanas: tiene que extraer el DOM o el texto del navegador, luego razonar, luego ejecutar clics o entradas de texto, y durante ese proceso pedirle información adicional o confirmación al ser humano; esto es la manifestación de que «el modo es la capacidad».
\end{importantbox}

\subsection{Planificación a largo plazo y limitaciones de observación}
El punto final del Web Agent casi nunca está en la página actual, por lo que debe mantener una estrategia interna para evitar descubrir, tras varios clics, que tomó un camino equivocado. Además, el SRT de 00:33:54 a 00:34:14 señala que, incluso cuando el Agent elabora un plan, también debe garantizar la precisión de las acciones que ejecuta, o de lo contrario tenderá a repetir pasos.

\begin{warningbox}{El riesgo de separar Planning y Execution}
El Web Agent puede elaborar un plan excelente, pero fallar porque la interfaz cambió, el texto del botón es ligeramente distinto o se requiere una verificación de varios niveles; en ese caso hay que retroalimentar el error hacia el tape, para que los nodos posteriores aprendan una manera de ejecución más robusta.
\end{warningbox}

\subsection{Resumen de la sección}
El Web Agent es el resultado combinado de que «el humano esté en el loop» y de que «el Agent esté en el loop»; la semántica de la interfaz de usuario, la planificación a largo plazo y la precisión en la ejecución de las acciones necesitan mecanismos sistemáticos que las garanticen.

